serves in the second-quantization formalism as the operator for the total
number of particles in the system. Indeed, on substituting the
$\psi$-operators from (64.20) and using the normalization and mutual
orthogonality of the wave functions, we find
$\hat N=\sum_i\hat a_i^+\hat a_i$. Each term in this sum is the particle-number
operator for the $i$th state; by (64.9), its eigenvalues are the occupation
numbers $N_i$. Their sum is the total number of particles in the
system\footnote{For systems having a prescribed number of particles, these
statements (as well as the properties of the free-particle Hamiltonian
(64.19)) appear trivial. In relativistic theory, however, their generalization
leads to new results which are far from trivial (cf. Vol.~IV, \S~11).}.
